% Drawing 43/03-3 -- folder 43, batch 03, archive page 59 (PDF page 60).
% Transcription: à droite, un croquis : un segment de $a$ à $b$, coupé entre $\alpha$ et $\beta$ par un trait oblique légendé « \uncertain{fenêtre} » ; dessous :
% Kind: segment a..b cut by an oblique stroke. Method 2 (schematic redraw from the facsimile). Status: ok.
% Notes: A segment from a to b, heavy from a to alpha, then from beta to b with a point and a tick; an oblique stroke cuts it between alpha and beta, labelled "fenetre" (uncertain reading kept).
% Coordinates: 1 unit = 1 mm = 1/100 of the crop width (bbox in page fractions: [0.6, 0.68, 0.98, 0.8]).
% Redrawn by Claude (claude-opus-5-5), issue #73, 2026-11-07; not checked by a human; not compiled (no TeX in the session).
% Rights: reproduces part of the work -- behind RENDER_TRANSCRIPTS, never published by itself.
\documentclass[tikz,border=3pt]{standalone}
\usepackage{amsmath,amssymb}
\usetikzlibrary{patterns,arrows.meta}
\tikzset{pencil/.style={black!45}, struck/.style={black!35}, mark/.style={circle,draw,inner sep=1.3pt}, pt/.style={circle,fill,inner sep=0.9pt}, opt/.style={circle,draw,fill=white,inner sep=1.1pt}}
\begin{document}
\begin{tikzpicture}[x=1mm,y=1mm,line cap=round,line join=round,font=\scriptsize]
\draw[very thick] (21,-17) -- (50,-17.5);
\draw[very thick] (21,-15.5) -- (21,-18.5);
\draw (53,-18) -- (93,-19);
\node[pt] at (57,-18) {};
\draw (64,-17) -- (64,-19.5);
\draw (93,-17.5) -- (93,-20.5);
\draw (61,-10) -- (54,-27);
\node[inner sep=1pt] at (21,-22) {\scriptsize $a$};
\node[inner sep=1pt] at (50,-22) {\scriptsize $\alpha$};
\node[inner sep=1pt] at (58,-24) {\scriptsize $\beta$};
\node[inner sep=1pt] at (92,-25) {\scriptsize $b$};
\node[inner sep=1pt] at (65,-7) {\tiny fenêtre};
\end{tikzpicture}
\end{document}
